\documentclass{amsart}
\usepackage{amsmath, amssymb, enumitem}
\title{Math 21b --- Course Learning Goals}
\date{}
\begin{document}
\maketitle

By the end of this course, students will be able to:

\bigskip

\textbf{Unit 1: Linear Systems and Matrix Algebra}
\begin{enumerate}[leftmargin=2em]
    \item Row-reduce a matrix to echelon form and RREF, and interpret the result in terms of solutions to a linear system.
    \item Determine whether a linear system has zero, one, or infinitely many solutions, and express the solution set in parametric form.
    \item Perform matrix addition, scalar multiplication, and matrix multiplication; compute transposes and inverses.
    \item Apply the Invertible Matrix Theorem to characterize invertible matrices in multiple equivalent ways.
\end{enumerate}

\bigskip

\textbf{Unit 2: Vector Spaces and Linear Transformations}
\begin{enumerate}[leftmargin=2em, resume]
    \item Define and verify that a set with operations forms a vector space or subspace.
    \item Find a basis and dimension for subspaces including null spaces, column spaces, and row spaces.
    \item State and apply the Rank-Nullity Theorem.
    \item Represent linear transformations as matrices and interpret kernel and image geometrically.
    \item Perform coordinate changes and change-of-basis computations.
\end{enumerate}

\bigskip

\textbf{Unit 3: Orthogonality and Least Squares}
\begin{enumerate}[leftmargin=2em, resume]
    \item Compute dot products, lengths, angles, and orthogonal projections.
    \item Apply the Gram--Schmidt process to produce an orthonormal basis.
    \item Solve least-squares problems by setting up and solving normal equations.
    \item Interpret the least-squares solution geometrically as an orthogonal projection.
\end{enumerate}

\bigskip

\textbf{Unit 4: Eigenvalues, Eigenvectors, and Diagonalization}
\begin{enumerate}[leftmargin=2em, resume]
    \item Find eigenvalues via the characteristic polynomial and compute corresponding eigenspaces.
    \item Determine when a matrix is diagonalizable and carry out the diagonalization $A = PDP^{-1}$.
    \item Use diagonalization to compute matrix powers and solve discrete dynamical systems.
    \item Analyze long-run behavior of dynamical systems using dominant eigenvalues.
\end{enumerate}

\bigskip

\textbf{Unit 5: Differential Equations and Fourier Series}
\begin{enumerate}[leftmargin=2em, resume]
    \item Solve first-order linear systems of differential equations using eigenvalue methods.
    \item Describe and classify solution behavior (stable, unstable, center) for linear systems.
    \item Compute Fourier coefficients and write the Fourier series of a piecewise-continuous periodic function.
    \item Apply Fourier series to analyze signals and solve basic partial differential equations (e.g., heat equation).
    \item Articulate connections between Fourier series and orthogonal projections in function spaces.
\end{enumerate}

\bigskip

\textbf{Overarching Goals}
\begin{enumerate}[leftmargin=2em, resume]
    \item Move fluidly between concrete (computational) and abstract (conceptual) perspectives on linear algebra.
    \item Recognize and model phenomena from the natural and social sciences using linear algebraic tools.
    \item Write clear, rigorous mathematical arguments and proofs.
    \item Communicate mathematical ideas precisely in both written and oral form.
\end{enumerate}

\end{document}
